\section{A unit exponent}
\label{sec:unit}

The complement quotient replaces coordinate symmetry when the other
two exponents are arbitrary. A root between zero and one then settles the
entire unit-exponent family after the small exceptions are identified.

\begin{theorem}
\label{thm:unit}
For distinct primes $p,q,r$ and every $b,c\geq1$, the graph
$G_{p q^b r^c}$ is not Laplacian integral. The same holds for
permutations of the exponent vector $(1,b,c)$.
\end{theorem}

The equal-exponent decomposition is replaced by a complement identity
on all six support classes. For arbitrary $a,b,c\geq1$, remove the
$u=abc-1$ universal vertices and denote the remaining graph by $H$.
In the order $1,2,3,12,13,23$, its class-size vector is
$w=(a,b,c,ab,ac,bc)^T$ and its vertex count is
$N=a+b+c+ab+ac+bc$. Disjoint supports are adjacent in the complement
$\overline H$, whose cell-constant Laplacian quotient is
\[
C=\begin{pmatrix}
b+c+bc&-b&-c&0&0&-bc\\
-a&a+c+ac&-c&0&-ac&0\\
-a&-b&a+b+ab&-ab&0&0\\
0&0&-c&c&0&0\\
0&-b&0&0&b&0\\
-a&0&0&0&0&a
\end{pmatrix}.
\]
Multiplication on the left by $\operatorname{diag}(w)$ makes $C$
symmetric. Moreover $C\mathbf1=0$ and $w^TC=0$.
If $B$ is the quotient of $H$, the complement relation gives
\begin{equation}
\label{eq:complement}
B+C=NI-\mathbf1w^T.
\end{equation}
Thus an eigenvector $Cv=\mu v$ with $\mu\ne0$ satisfies $w^Tv=0$
and $Bv=(N-\mu)v$. When $N-\mu\ne0$, Lemma~\ref{lem:join}
lifts it to an eigenvalue $N+u-\mu$ of $G_n$.

\begin{proof}[Proof of Theorem~\ref{thm:unit}]
Write $\det(xI-C)=x h_{a,b,c}(x)$. Direct determinant expansion yields
\begin{align*}
h_{a,b,c}(0)&=-abc(a+b+c)N,\\
h_{a,b,c}(a)&=-a^2b^2c^2(a^2+2a-b-c).
\end{align*}
At $a=1$, the first value is strictly negative and
\[
h_{1,b,c}(1)=b^2c^2(b+c-3).
\]
When $b+c>3$, a real root $\mu$ lies strictly between $0$ and $1$.
Equation~\eqref{eq:complement} and the universal-vertex lift produce a
Laplacian eigenvalue $V-\mu\in(V-1,V)$, where $V=N+u$ is the integer
vertex count. This eigenvalue is noninteger. When $b+c\leq3$, the
positive pairs are $(1,1),(1,2),(2,1)$; Theorem~\ref{thm:pqrk}
settles the corresponding permutations of $(1,1,k)$. Permuting the
prime factors preserves the graph's isomorphism type.
\end{proof}

The unit-exponent obstruction extends to any number of prime factors
by a variational argument on the complement graph.

\begin{theorem}
\label{thm:unit-any-prime-count}
Let $n=\prod_{i=1}^t p_i^{k_i}$ with $t\geq3$, distinct primes
$p_i$ and positive integer exponents $k_i$. If some $k_i=1$, then
$G_n$ has a Laplacian eigenvalue in $(V-1,V)$, where
$V=\prod_{i=1}^t(k_i+1)-2$. In particular, $G_n$ is not
Laplacian integral.
\end{theorem}
\begin{proof}
Permute the factors so that $k_t=1$, and put
\[
 A=\prod_{i<t}k_i,\qquad
 B=\prod_{i<t}(k_i+1)-1-A.
\]
The quantity $B$ is the sum of the weights of all nonempty proper
supports on the first $t-1$ coordinates. Since $t-1\geq2$, it is
positive. Remove the $A-1$ full-support universal vertices and let
$F=\overline H$. Then $N=|H|=A+2B+1$ and $V=N+A-1=2A+2B$.
The graph $F$ is connected: its singleton support classes form a
connected complete multipartite graph, and every proper support has
a singleton outside it.

The support $\{t\}$ consists of a single vertex $z$. Each of the
$A$ vertices supported on $[t-1]$ has $z$ as its sole neighbor in
$F$. The remaining $2B$ vertices form two collections of weight $B$:
nonempty proper supports $S\subset[t-1]$, and their counterparts
$S\cup\{t\}$. The vertex $z$ is adjacent to all vertices of the
first collection and none of the second.

Give a function $f$ the value $2B$ on the $A$ leaves, zero at $z$,
and $-A$ on the other $2B$ vertices. It is nonzero and
orthogonal to the constant vector. Its Laplacian energy and norm are
\[
 f^TL(F)f=A(2B)^2+BA^2=AB(4B+A),\qquad
 f^Tf=A(2B)^2+2BA^2=2AB(2B+A).
\]
Indeed, only edges incident to $z$ have unequal endpoint values.
Connectedness and the variational characterization therefore give
\[
 0<\lambda_2(F)\leq\frac{4B+A}{4B+2A}<1.
\]
For a corresponding eigenvector, the complement identity
$L(H)+L(F)=NI-J$ gives the nonzero eigenvalue $N-\lambda_2(F)$
of $H$. Lemma~\ref{lem:join} lifts it to
$V-\lambda_2(F)\in(V-1,V)$, proving the claim.
\end{proof}

The restriction $t\geq3$ is sharp for this conclusion. If $t=2$
and $k_2=1$, then $B=0$ and $F=K_{1,A}$, so the complement and
universal-vertex identities give an integral spectrum for $G_n$.
Theorem~\ref{thm:unit-any-prime-count} also gives a proof of
Theorems~\ref{thm:squarefree} for $t\geq3$ and~\ref{thm:unit}
without their separate spectral decompositions or small exceptions.

The general six-support reduction also gives an exact diagnostic for
fully distinct triples. For $1\leq a<b<c\leq7$, all thirty-five quintics
have a nonlinear irreducible factor over $\mathbb Q$. The quintics
at $(1,2,7)$ and $(2,3,5)$ have factor degrees $(1,4)$, and that at
$(2,4,6)$ has factor degrees $(2,3)$;
irreducibility of the entire quintic is therefore not a
uniform obstruction. Theorem~\ref{thm:unit} covers every triple beginning
with $1$ without a range restriction; the thirty-five-case calculation
itself remains finite evidence.
